\section{An edge is a contract}
\label{sec:gredge}

What does a desk actually believe when it draws an edge?  Listen to
the sentences.  ``The sister moves nine tenths of a point when NVDA
moves a point.''  ``A six-month move of one point is a two-point move
at three months.''  Each sentence has two independent clauses hiding
in it: \emph{how large} the implied move is, and \emph{how much the
speaker trusts the relation}.  A trader can assert a large amplitude
with little confidence, or a modest amplitude with great confidence.
Any encoding that forces the two into one number will betray one of
them.

So an edge carries two numbers.  The relation ``node $j$ informs node
$i$'' is the statement
\begin{equation}
\Theta_i \;=\; \beta_{ij}\,\Theta_j \;+\; \text{noise},
\qquad
\text{noise mean }0,\ \text{variance } 1/p_{ij}.
\label{eq:gredge}
\end{equation}
The \emph{amplitude} $\beta_{ij}$ says what the receiver's move should
be per unit of the informer's move --- it is the finance ``beta,''
deliberately: the move-per-move slope a desk quotes between related
assets --- and it prices the \emph{mean}.  The \emph{relation
precision} $p_{ij}$ says how tightly the relation is believed ---
precision is one over variance, as everywhere since
\cref{sec:infgate} --- and it prices the \emph{band}.  Amplitude and
trust never mix: that is the contract, and this section and the next
make it hold exactly, not approximately.

\subsection{Where a calendar amplitude comes from}

Cross-name amplitudes are desk judgments (a name's beta to its
index).  Calendar amplitudes --- between two expiries of one name ---
can be \emph{derived}, and the derivation is two lines.  Recall the
chart identity in force since Chapter~2: a slice's total implied
variance is $w=\sigma^{2}\tau$, volatility squared times variance
time.  Suppose a piece of news adds the same amount of total variance
$\delta w$ to every expiry --- the natural first model of a systematic
repricing: the market now expects one more shock's worth of variance,
and every expiry that contains the shock inherits all of it (this is
Chapter~8's event logic wearing scenario clothes).  Differentiate
$w=\sigma^{2}\tau$ at fixed $\tau$:
\[
\delta w \;=\; 2\,\sigma\,\tau\,\delta\sigma
\qquad\Longrightarrow\qquad
\delta\sigma \;=\; \frac{\delta w}{2\,\sigma\,\tau}.
\]
The same $\delta w$ lands on every expiry, so the \emph{volatility}
moves scale as
\[
\frac{\delta\sigma_i}{\delta\sigma_j}
=\frac{\sigma_j\,\tau_j}{\sigma_i\,\tau_i}
\;\approx\;\frac{\tau_j}{\tau_i},
\]
the approximation being a flat volatility term structure
($\sigma_i\approx\sigma_j$).  The chapter's calendar contracts adopt
the clean form as the stated convention:
\begin{equation}
\beta_{ij}=\frac{\tau_j}{\tau_i}
\qquad\text{(calendar edges; constant total-variance injection)}.
\label{eq:grcalbeta}
\end{equation}
Read it on a clean ladder of three expiries with $\tau=0.25$, $0.5$
and $1.0$ (variance time, as always) --- call the nodes 3M, 6M and 1Y.  A $+1$-point move at 6M implies
$+1\times(0.5/0.25)=+2$ points at 3M and $+1\times(0.5/1.0)=+0.5$ at
1Y.  Short-dated volatility swings hardest --- the same $\delta w$ is a
bigger share of a smaller $w$ --- which is precisely what
\cref{fig:gruniverse}(b) showed on the real board.  Note also the
built-in consistency: $\beta_{ij}\beta_{ji}
=(\tau_j/\tau_i)(\tau_i/\tau_j)=1$, so a ladder cannot claim
amplification in both directions, and around any cycle of calendar
edges the amplitudes telescope to one.  (Asserted cross-name
amplitudes carry no such guarantee; a loop of them that fails to
multiply to one asserts contradictory scalings, and the solve of the
next section will then split the difference somewhere the desk did not
choose.  The reference implementation flags such loops rather than
solving through them silently.)

\subsection{One relation, two readings}

Statement \cref{eq:gredge} is written with $i$ as receiver, but a
relation between two nodes is one fact, not two.  Rewrite the same
noise term in the other node's units: if
$\Theta_i-\beta\,\Theta_j$ has variance $1/p$, then dividing by
$\beta$ shows that
\[
\Theta_j-\frac{1}{\beta}\,\Theta_i
\quad\text{has variance}\quad
\frac{1}{\beta^{2}}\cdot\frac{1}{p}
\qquad\Longrightarrow\qquad
p_{\,\text{reverse}}=\beta^{2}\,p .
\]
The reverse reading of the contract carries the reciprocal amplitude
$1/\beta$ --- as it must --- and a precision rescaled by $\beta^{2}$,
because the same uncertainty is being quoted in different units.  This
is bookkeeping, not a second belief.  The convention that keeps it
tidy (the \emph{canonical orientation}) is to quote each calendar
relation with the \emph{shorter} maturity as receiver: relation noise
is then stated in short-dated volatility points, where moves are
largest and desk intuition lives.  One consequence is worth saying
now, because \cref{sec:grbelieve} will lean on it: a contract is
symmetric information.  If the relation between $A$ and $B$ is
trusted tightly, it is trusted tightly \emph{in both directions} ---
observing $B$ legitimately updates $A$.  The arrow in a diagram
chooses units and bookkeeping; it does not create one-way inference.
Making influence genuinely one-way is a different, stronger assertion,
and it costs machinery this section does not have.

\subsection{The contract, tested}

\Cref{fig:grcontract} runs the two-dial claim through the chapter's
solver on the clean ladder, 6M lit at $+1$ point, sweeping the
relation precision across three decades.  In panel~(a), the posterior
means at 3M and 1Y sit at exactly $+2.00$ and $+0.50$ --- the
configured amplitudes --- across the entire sweep; the largest
deviation anywhere is \MacGrContractFlat\ points, solver arithmetic
rather than economics.  Turning the trust dial by a factor of a
thousand leaves the answer where the amplitude put it: amplitude is
not a function of confidence.  Panel~(b) shows
what the trust dial \emph{does} govern: the posterior standard
deviations fall exactly as $1/\sqrt{p}$ (the dotted reference line lies
on the 3M curve).  And the two curves sit a fixed factor apart --- the
measured 3M/1Y ratio is \MacGrContractSdRatio\ --- which is the units
law just derived, in two steps the reader can retrace.  The 6M--1Y
relation is quoted with 6M as receiver at amplitude $2$, noise
variance $1/p$; its reverse reading, the one the 1Y node hears,
carries precision $\beta^{2}p=4p$, hence standard deviation
$\tfrac12\cdot1/\sqrt{p}$ --- half the 3M node's $1/\sqrt{p}$.

\begin{figure}[htbp]
\centering
\paperfig{\textwidth}{fig_gr_contract.pdf}
\caption{The contract on the clean ladder (3M/6M/1Y, 6M lit at $+1$
volatility point, calendar amplitudes \cref{eq:grcalbeta}).
(a)~Posterior means at the two dark nodes against the relation
precision: flat at the configured $+2.00$ and $+0.50$ across three
decades (maximum deviation \MacGrContractFlat).  (b)~Posterior
standard deviations: falling as $1/\sqrt{p}$ (dotted), with the 1Y
band exactly half the 3M band --- the same relation noise read in
1Y units.  Amplitude prices the mean; trust prices the band.}
\label{fig:grcontract}
\end{figure}

We now know what one edge says.  A universe, though, gives every node
several edges, and several lit nodes speak at once.  The next section
lets the messages meet.
